\documentclass[11pt,a4paper]{article}
\usepackage{fontspec}
\usepackage{polyglossia}
\setdefaultlanguage{russian}
\setmainfont{PT Serif}
\setsansfont{PT Sans}
\setmonofont{PT Mono}
\newfontfamily\cyrillicfont{PT Serif}
\newfontfamily\cyrillicfonttt{PT Mono}
\usepackage[margin=25mm]{geometry}
\usepackage{tabularx,longtable,fancyvrb}
\usepackage[hidelinks]{hyperref}
\setlength{\parindent}{1.25em}
\setlength{\emergencystretch}{3em}
\setlength{\parskip}{0.3em}
\renewcommand{\arraystretch}{1.25}
\hypersetup{pdftitle={ИЫ.00110-01 30 01},pdfsubject={ВЫЧИСЛИТЕЛЬНЫЙ КОМПЛЕКС "ЭЛЬБРУС 1-КБ" ОБЩЕЕ СИСТЕМНОЕ ПРОГРАММНОЕ ОБЕСПЕЧЕНИЕ ФОРМУЛЯР}}
\begin{document}
\begin{titlepage}\small
\noindent\begin{minipage}[t]{0.45\linewidth}Согласовано\\Представитель заказчика №\,2272\\В.~И.~Пономаренко\end{minipage}\hfill\begin{minipage}[t]{0.45\linewidth}\raggedleft Утверждаю\\Руководитель предприятия\\Г.~Г.~Рябов\end{minipage}
\vspace*{\fill}\begin{center}
{\large Вычислительный комплекс «Эльбрус 1-КБ»}\\[0.6em]
{\large Общее системное программное обеспечение}\\[0.6em]
{\large Формуляр}\\[0.6em]
\vspace{1em}{\large\bfseries ИЫ.00110-01 30 01}\\[0.4em]
Листов 30\end{center}\vspace*{\fill}
\begin{flushright}\begin{tabular}{r}Представители предприятия-разработчика\\Начальник отдела\\В.~Ф.~Тюрин\\Исполнители\\Е.~В.~Голенкова\\А.~А.~Нахипова\\Нормоконтролер\\Т.~П.~Рудаченко\end{tabular}\end{flushright}
\begin{center}1988\end{center}\end{titlepage}
\tableofcontents\clearpage
\section*{1.~Общие указания}\addcontentsline{toc}{section}{1.~Общие указания}
\par\textbf{1.1.} Настоящий формуляр распространяется на программное изделие «Общее системное программное обеспечение» ИЫ.00110-01 (ОСПО) ВК «Эльбрус 1-КБ» и является документом, удостоверяющим гарантированные предприятием-изготовителем основные характеристики ОСПО ВК «Эльбрус 1-КБ», определяющим комплект поставки и содержащим сведения об изменениях в эксплуатационной документации и текстах программ ОСПО ВК «Эльбрус 1-КБ» на магнитных носителях данных (МНД).\par
\par\textbf{1.2.} Формуляр должен находиться в подразделении, ответственном за эксплуатацию общего системного программного обеспечения.\par
\par\textbf{1.3.} Все записи в формуляре должны производиться чернилами (пастой) или тушью, отчетливо и аккуратно. Подчистки, помарки и незаверенные исправления не допускаются.\par
\par\textbf{1.4.} Разделы 4, 6, 7 заполняются предприятием --- Изготовителем, а разделы 5, 9, 10, 11, 12, 13 --- Пользователем ВК «Эльбрус 1-КБ».\par
\par\textbf{1.5.} Получение рабочей копии ОСПО ВК «Эльбрус 1-КБ» с архивной копии производится в соответствии с приложением 1 формуляра.\par
\par\textbf{1.6.} Перед эксплуатацией ОСПО необходимо ознакомиться со всеми документами, перечисленными в ведомости эксплуатационных документов ИЫ.00110-01 20 01. Получение документов на бумажном носителе производится в соответствии с приложением 2 формуляра.\par
\section*{2.~Общие сведения}\addcontentsline{toc}{section}{2.~Общие сведения}
\par\textbf{2.1.} Общее системное программное обеспечение (ОСПО) ИЫ.00110-01 предназначено для организации решения широкого круга научных, производственных и специальных задач в режиме мультипрограммирования на вычислительном комплексе ВК «Эльбрус 1-КБ».\par
\par\textbf{2.2.} ОСПО ВК «Эльбрус 1-КБ» поставляется пользователю в соответствии с разделом 4 настоящего документа.\par
\section*{3.~Основные характеристики}\addcontentsline{toc}{section}{3.~Основные характеристики}
\par\textbf{3.1.} ОСПО ИЫ.00110-01 обеспечивает работу пользователей на вычислительном комплексе ВК «Эльбрус 1-КБ» в пакетном, диалоговом и диалогово-пакетном режимах как на автокодах (автокоды ИТМ, БЕМШ, MADLEN), так и на языках высокого уровня (Алгол, Фортран, Паскаль).\par
\par\textbf{3.2.} В состав ОСПО ВК «Эльбрус 1-КБ» входят следующие программные компоненты:\par
\par\hangindent=1.5em - Операционная система Диспак;\par
\par\hangindent=1.5em - Системы программирования режима совместимости;\par
\par\hangindent=1.5em - Система программирования, включающая ассемблер\par
ВК «Эльбрус 1-КБ».

\par\textbf{3.3.} ОСПО ВК «Эльбрус 1-КБ» обеспечивает:\par
\par\hangindent=1.5em 1) Непрерывную мультипрограммную работу вычислительной системы в режимах пакетной обработки, дистанционно-пакетной обработки, работу пользователей с терминалов в режиме разделения времени, совмещение указанных режимов;\par
\par\hangindent=1.5em 2) Формирование и накопление задач во входной очереди до 255;\par
\par\hangindent=1.5em 3) Возможность управления планированием мультипрограммной смеси задач в соответствии с потребностями вычислительного центра;\par
\par\hangindent=1.5em 4) Ведение и контроль бюджетно-статистического учета использования ресурсов;\par
\par\hangindent=1.5em 5) Средства бюджетирования архивного пространства на магнитных дисках и защиту информации пользователя от несанкционированного доступа.\par
\par\textbf{3.4.} Характеристиками общего системного программного обеспечения, удостоверяющими правильность записи, считывания и хранения на магнитных лентах, являются контрольные суммы магнитных лент (см. ТАБЛ.1), на которых поставляется ОСПО.\par
\begin{flushright}\small Таблица 1\end{flushright}
\begin{center}\small
\begin{BVerbatim}
-----------------------------------------
     Номер МЛ     I      контрольная
                  I         сумма
------------------I----------------------
                  I
       МЛ1        I  7205 7233 3054 0257
                  I
       МЛ2        I  5224 5163 0300 7443
                  I
       МЛ3        I  5750 0101 7656 6260
                  I
       МЛ4        I  0035 7507 0150 4344
                  I
                  I
\end{BVerbatim}
\end{center}

\section*{4.~Комплектность}\addcontentsline{toc}{section}{4.~Комплектность}
ОСПО должно предъявляться комплектно согласно ведомости эксплуатационных документов ИЫ.00110-01 20 01 и настоящему формуляру.

При этом титульные части, информационные части и части регистрации изменений всех документов поставляются на бумаге, основные части всех документов поставляются на магнитном носителе. Исключение составляют документы «Ведомость эксплуатационных документов» ИЫ.00110-01 20 01 и «Формуляр» ИЫ.00110-01 30 01, которые поставляются полностью на бумаге.

\begin{center}\small\begin{tabularx}{\linewidth}{|>{\raggedright\arraybackslash\hsize=1.418\hsize}X|>{\raggedright\arraybackslash\hsize=1.269\hsize}X|>{\raggedright\arraybackslash\hsize=0.522\hsize}X|>{\raggedright\arraybackslash\hsize=0.821\hsize}X|>{\raggedright\arraybackslash\hsize=0.970\hsize}X|}\hline
\small\bfseries\hspace{0pt}Обозначение & \small\bfseries\hspace{0pt}Наименование & \small\bfseries\hspace{0pt}Количество & \small\bfseries\hspace{0pt}Порядковый учетный номер & \small\bfseries\hspace{0pt}Примечание \\ \hline
 & Документы &  &  &  \\ \hline
ИЫ.00110-01 20 01 & Комплект эксплуатационной документации согласно ведомости эд магнитные носители & 1 &  & На бумаге \\ \hline
МЛ1 МЛ2 МЛ3 МЛ4 & ГП ОС Диспак системы программирования режима совместимости система программирования, эксплуатационная документация ОС Диспак. Текст программы (файлы 1-2); система программирования. Текст программы (файлы 3-8); система программирования. ассемблер. Текст программы (файлы 9-14) & 1 1 1 1 &  & Не поставляется \\ \hline
\end{tabularx}\end{center}

\section*{5.~Периодический контроль основных характеристик}\addcontentsline{toc}{section}{5.~Периодический контроль основных характеристик}
\begin{center}При эксплуатации и хранении\end{center}
Контроль работоспособности ОСПО осуществляется один раз в год путем подсчета контрольных сумм каждой магнитной ленты и сличения этих сумм с эталонными значениями, приведенными в ТАБЛ.1 настоящего формуляра. Результаты проверки заносятся в ТАБЛ.2 настоящего формуляра. В графе «Величина» проставляются контрольные суммы МЛ копии ОСПО, находящейся у пользователя ВК «Эльбрус 1-КБ» к моменту периодического контроля. Способ подсчета контрольных сумм приведен в документе «Операционная система Диспак. Руководство системного программиста» ИЫ.00110-01 32 01, раздел 3.

\begin{flushright}\small Таблица 2\end{flushright}
\begin{center}\small
\begin{BVerbatim}
─────────────────────────────────────────────────────────────────
   Проверяемая   │         Дата   проведения   измерения
  характеристика │───────────────────────────────────────────────
─────────────────│    19   г.    │    19   г.    │    19   г.
 наиме-│         │───────────────────────────────────────────────
 нова- │         │ факти-│замерил│ факти-│замерил│ факти-│замерил
 ние   │  вели-  │ ческая│(долж- │ ческая│(долж- │ ческая│(долж-
 изме- │  чина   │ вели- │ ность,│ вели- │ ность,│ вели- │ ность,
 рения │         │ чина  │ под-  │ чина  │ под-  │ чина  │ под-
       │         │       │ пись) │       │ пись) │       │ пись)
 ──────│─────────│───────│───────│───────│───────│───────│───────
       │         │       │       │       │       │       │
       │         │       │       │       │       │       │
       │         │       │       │       │       │       │
       │         │       │       │       │       │       │
       │         │       │       │       │       │       │
       │         │       │       │       │       │       │
       │         │       │       │       │       │       │
       │         │       │       │       │       │       │
       │         │       │       │       │       │       │
       │         │       │       │       │       │       │
       │         │       │       │       │       │       │
       │         │       │       │       │       │       │
       │         │       │       │       │       │       │
       │         │       │       │       │       │       │
       │         │       │       │       │       │       │
       │         │       │       │       │       │       │
       │         │       │       │       │       │       │
       │         │       │       │       │       │       │
       │         │       │       │       │       │       │
       │         │       │       │       │       │       │
       │         │       │       │       │       │       │
       │         │       │       │       │       │       │
       │         │       │       │       │       │       │
       │         │       │       │       │       │       │
       │         │       │       │       │       │       │
       │         │       │       │       │       │       │
       │         │       │       │       │       │       │
       │         │       │       │       │       │       │
       │         │       │       │       │       │       │
       │         │       │       │       │       │       │
       │         │       │       │       │       │       │
       │         │       │       │       │       │       │
       │         │       │       │       │       │       │
       │         │       │       │       │       │       │
       │         │       │       │       │       │       │
       │         │       │       │       │       │       │
       │         │       │       │       │       │       │
       │         │       │       │       │       │       │
       │         │       │       │       │       │       │
       │         │       │       │       │       │       │
       │         │       │       │       │       │       │
\end{BVerbatim}
\end{center}

\section*{6.~Свидетельство о приемке}\addcontentsline{toc}{section}{6.~Свидетельство о приемке}
\subsection*{6.1.~Общее системное программное обеспечение ИЫ.00110-01}\addcontentsline{toc}{subsection}{6.1.~Общее системное программное обеспечение ИЫ.00110-01}
\begin{center}\small
\begin{BVerbatim}
---------------------------------------    -----------
\end{BVerbatim}
\end{center}

Соответствует техническим условиям ИЫ.00110-01 28 01

\begin{center}\small
\begin{BVerbatim}
                                         -----------------
И признано  годным  для  эксплуатации.
                              Дата выпуска
\end{BVerbatim}
\end{center}

\begin{center}\small
\begin{BVerbatim}
                                               -------------
Руководитель предприятия               начальник ОТК
    М.П.
\end{BVerbatim}
\end{center}

\subsection*{6.2.~Заключение представителя заказчика}\addcontentsline{toc}{subsection}{6.2.~Заключение представителя заказчика}
\begin{center}М.П. представитель заказчика\end{center}
\begin{center}\small
\begin{BVerbatim}
-----------------------
      «  »      19  г.
          ------
\end{BVerbatim}
\end{center}

7. Свидетельство об упаковке и маркировке изделие ОСПО ИЫ.00110-01

\begin{center}\small
\begin{BVerbatim}
    ------------                              -----------
Номер программного изделия предприятия
                                      ---------------
Упакован(о)
           ------------------------------------------
Согласно требованиям, предусмотренным
Инструкцией
           ------------------------------------------
      Дата упаковки
                   ---------------
      Упаковку произвел                        (подпись)
                       ----------------------
      Изделие после упаковки принял                 (подпись)
                                   ---------------
                                                      М.П.
\end{BVerbatim}
\end{center}

\section*{8.~Гарантийные обязательства}\addcontentsline{toc}{section}{8.~Гарантийные обязательства}
Предприятие-изготовитель гарантирует характеристики, приведенные в разделе 3, а также сохранность и воспроизведение информации с поставленных магнитных носителей при хранении копии ОСПО ВК «Эльбрус 1-КБ» в условиях, определенных

В разделе 10.

Срок гарантии определяется сроком действия гарантийных обязательств на ВК «Эльбрус 1-КБ».

\section*{9.~Сведения о рекламациях}\addcontentsline{toc}{section}{9.~Сведения о рекламациях}
\begin{center}\small\begin{tabularx}{\linewidth}{|>{\raggedright\arraybackslash\hsize=0.424\hsize}X|>{\raggedright\arraybackslash\hsize=1.030\hsize}X|>{\raggedright\arraybackslash\hsize=1.394\hsize}X|>{\raggedright\arraybackslash\hsize=1.152\hsize}X|}\hline
\small\bfseries\hspace{0pt}Дата & \small\bfseries\hspace{0pt}Содержание рекламации & \small\bfseries\hspace{0pt}Меры, принятые по рекламации & \small\bfseries\hspace{0pt}Подпись ответственного лица \\ \hline
\end{tabularx}\end{center}

\section*{10.~Сведения о хранении}\addcontentsline{toc}{section}{10.~Сведения о хранении}
Хранение копии ОСПО ВК «Эльбрус 1-КБ» должно производиться в помещении при температуре от +5 до +40 градусов цельсия и относительной влажности воздуха от 30\% до 95\%.

Нельзя хранить магнитные ленты вблизи сильных электромагнитных полей (электромагнитов, электромоторов, трансформаторов и Т.Д.).

Магнитные ленты должны храниться вертикально, вдали (не менее одного метра) от нагревательных приборов и должны быть защищены от воздействия прямых солнечных лучей.

\begin{center}\small
\begin{BVerbatim}
─────────────────────────────────────────────────────────────────
          Дата          │          │
────────────────────────│          │     должность, фамилия и
 установки  │ снятия с  │ условия  │ подпись лица, ответственного
на хранение │ хранения  │ хранения │       за хранение
────────────│───────────│──────────│─────────────────────────────
            │           │          │
            │           │          │
            │           │          │
            │           │          │
            │           │          │
            │           │          │
            │           │          │
            │           │          │
            │           │          │
            │           │          │
            │           │          │
            │           │          │
            │           │          │
            │           │          │
            │           │          │
            │           │          │
            │           │          │
            │           │          │
            │           │          │
            │           │          │
            │           │          │
            │           │          │
            │           │          │
            │           │          │
            │           │          │
            │           │          │
            │           │          │
            │           │          │
            │           │          │
            │           │          │
            │           │          │
            │           │          │
            │           │          │
            │           │          │
            │           │          │
            │           │          │
            │           │          │
            │           │          │
            │           │          │
            │           │          │
            │           │          │
            │           │          │
            │           │          │
            │           │          │
            │           │          │
            │           │          │
            │           │          │
            │           │          │
\end{BVerbatim}
\end{center}

\section*{11.~Сведения о закреплении программного изделия}\addcontentsline{toc}{section}{11.~Сведения о закреплении программного изделия}
\begin{center}При эксплуатации\end{center}
\begin{center}\small
\begin{BVerbatim}
─────────────────────────────────────────────────────────────────
    Должность   │    Фамилия     │ Номер и дата приема │ Подпись
 ответственного │ ответственного │──────────────────── │ ответст-
     лица       │    лица        │ о назна- │ об осво- │ венного
                │                │  чении   │ бождении │  лица
 ───────────────│────────────────│──────────│──────────│─────────
                │                │          │          │
                │                │          │          │
                │                │          │          │
                │                │          │          │
                │                │          │          │
                │                │          │          │
                │                │          │          │
                │                │          │          │
                │                │          │          │
                │                │          │          │
                │                │          │          │
                │                │          │          │
                │                │          │          │
                │                │          │          │
                │                │          │          │
                │                │          │          │
                │                │          │          │
                │                │          │          │
                │                │          │          │
                │                │          │          │
                │                │          │          │
                │                │          │          │
                │                │          │          │
                │                │          │          │
                │                │          │          │
                │                │          │          │
                │                │          │          │
                │                │          │          │
                │                │          │          │
                │                │          │          │
                │                │          │          │
                │                │          │          │
                │                │          │          │
                │                │          │          │
                │                │          │          │
                │                │          │          │
                │                │          │          │
                │                │          │          │
                │                │          │          │
                │                │          │          │
                │                │          │          │
                │                │          │          │
                │                │          │          │
                │                │          │          │
                │                │          │          │
                │                │          │          │
\end{BVerbatim}
\end{center}

\section*{12.~Сведения об изменениях}\addcontentsline{toc}{section}{12.~Сведения об изменениях}
\begin{center}\small\begin{tabularx}{\linewidth}{|>{\raggedright\arraybackslash\hsize=1.727\hsize}X|>{\raggedright\arraybackslash\hsize=0.727\hsize}X|>{\raggedright\arraybackslash\hsize=0.727\hsize}X|>{\raggedright\arraybackslash\hsize=0.818\hsize}X|>{\raggedright\arraybackslash\hsize=1.091\hsize}X|>{\raggedright\arraybackslash\hsize=0.909\hsize}X|}\hline
\small\bfseries\hspace{0pt}Основание (входящий номер сопро-водительного документа и дата) & \small\bfseries\hspace{0pt}Дата проведения изменения & \small\bfseries\hspace{0pt}Содержание изменения & \small\bfseries\hspace{0pt}Поряд-ковый номер изменения & \small\bfseries\hspace{0pt}Должность, фамилия и подпись ответственного лица за проведение изменения & \small\bfseries\hspace{0pt}Подпись лица, ответственного за эксплуатацию программного изделия \\ \hline
\end{tabularx}\end{center}

\section*{13.~Особые отметки}\addcontentsline{toc}{section}{13.~Особые отметки}
\section*{Приложение 1.}\addcontentsline{toc}{section}{Приложение 1.}
\begin{center}Получение рабочей копии ОСПО\end{center}
Для получения рабочей копии ОСПО необходимо разметить магнитные диски с номерами 2053 и 2048 (СМ.п.~1 настоящего приложения), переписать информацию с МЛ1 на размеченный диск 2053 с помощью пультовой программы и вызвать операционную систему. (СМ.п.~2 настоящего приложения). На работающей операционной системе с помощью программы «UNICOP» переписать информацию с МЛ2 и МЛ3 на размеченный магнитный диск 2048 (СМ.п.~3 настоящего приложения). После окончания переписи на системных дисках 2053 и 2048 получена рабочая копия ОСПО. При переписи программой «UNICOP» на АЦПУ выдается контрольная сумма переписанной магнитной ленты.

\section*{1.~Разметка магнитных дисков}\addcontentsline{toc}{section}{1.~Разметка магнитных дисков}
Разметка магнитных дисков осуществляется с помощью автономной программы разметки (без ОС), см. «Пультовое программное обеспечение. Руководство оператора»

ИЫ.00112-01 34 01. Для ее вызова необходимо поставить на накопитель гибких магнитных дисков пультовой машины (ПМ) гибкий магнитный диск с программами разметки МД и вызова ОС. С помощью средств пультового программного обеспечения (ППО) набрать на клавиатуре пультовой машины следующую информацию (номера ячеек и содержимое регистров приведены в шестнадцатиричном виде):

\par\hangindent=1.5em - У0\par
\par\hangindent=1.5em - Заг SERWIS\par
Ответ системы: есть --- Для ПМ типа ДВК, начальный адрес программы и длина ее в словах --- Для ПМ типа ЕС-1840.

Задать параметры разметки:

\par\hangindent=1.5em - Ячейка 2:= 0\par
\par\hangindent=1.5em - Ячейка 3:= N пакета в шестнадцатиричном виде\par
\par\hangindent=1.5em - Ячейка в:= N устройства в шестнадцатиричном виде\par
\par\hangindent=1.5em - Ячейка с:= N канала в шестнадцатиричном виде\par
\par\hangindent=1.5em - У0\par
\par\hangindent=1.5em - Ар\par
\par\hangindent=1.5em - Убедиться, что РК=0 3F 07 000 (стоп 0)\par
\par\hangindent=1.5em - Директива «П»- пуск или нажатие соответствующей клавиши\par
На клавиатуре пультовой машины (ПМ) (таб --- Для ДВК-2М, НТ --- Для ДВК-3, печ --- Для ЕС-1840: см. «Пультовое программное обеспечение. Руководство оператора» ИЫ.00112-01 34 01).

В процессе разметки номер размечаемых цилиндра и головки заносятся в ячейку 13. После окончания разметки программа выходит на «Стоп 0» ( на РК код 0 3F 07 000).

\section*{2.~Постановка ОСПО на ВК «Эльбрус 1-КБ»}\addcontentsline{toc}{section}{2.~Постановка ОСПО на ВК «Эльбрус 1-КБ»}
Для постановки ОСПО (ОС и систем программирования) необходимо:

\par\hangindent=1.5em 1) Включить вычислительный комплекс в соответствии с инструкцией по эксплуатации ВК Э1-КБ (ПЭ1.700.018 иэ);\par
\par\hangindent=1.5em 2) Подготовить к работе пультовую машину и вызвать программу имитации инженерного пульта (ПП) (см. «Пультовое программное обеспечение. Руководство оператора». ИЫ.00112-01 34 01);\par
\par\hangindent=1.5em 3) Разметить полностью два диска с номерами 2053 и 2048 с помощью автономной программы разметки (СМ п.~1 настоящего приложения);\par
\par\hangindent=1.5em 4) Переписать готовые программы ОС с транспортной ленты МЛ1 на диск 2053, воспользовавшись для этого автономной программой переписи. Магнитная лента должна быть установлена на МГ 32 (на такой номер магнитофона настроена пультовая программа ПП13).\par
Для переписи необходимо выполнить следующую последовательность действий в рамках ПП:

У0 --- Установка нуля

ПП13 --- Чтение с МЛ в ОП

Директива «П»- пуск или нажатие соответствующей клавиши на ПМ;

Убедиться, что РК = 0 3F 07 000 (стоп 0);

Установить параметры чтения с МЛ:

Слово 5 --- N канала (стандартно = 3)

Слово 6 --- N устройства (стандартно = 2)

Слово 7 --- N файла (стандартно = 3)

Слово 8 --- 0 --- Читаемый код исполняется

\begin{center}В режиме совместимости с БЭСМ-6\end{center}
\begin{center}(Стандартно = 0)\end{center}
Директива «П»- пуск или нажатие соответствующей клавиши на ПМ;

Убедиться, что чтение закончено и РК = 0 3F 07 000;

Затем необходимо выполнить последовательность действий для переписи информации из ОП на МД:

\par\hangindent=1.5em - У0\par
\par\hangindent=1.5em - Заг SERWIS\par
Ответ системы: есть --- Для ПМ типа ДВК, начальный адрес программы и длина ее в словах --- Для ПМ типа ЕС-1840.

\par\hangindent=1.5em - Ячейка 2:= 1\par
\par\hangindent=1.5em - Ячейка 3:= N пакета в шестнадцатиричном виде\par
\par\hangindent=1.5em - Ячейка в:= N устройства МД в шестнадцатиричном виде\par
\par\hangindent=1.5em - Ячейка с:= N канала МД в шестнадцатиричном виде\par
\par\hangindent=1.5em - У0\par
\par\hangindent=1.5em - Ар\par
\par\hangindent=1.5em - Убедиться, что РК=0 3F 07 000 (стоп 0)\par
\par\hangindent=1.5em - Директива «П»- пуск или нажатие соответствующей клавиши на ПМ; после окончания переписи программа выходит на «Стоп 0» ( на РК код 0 3F 07 000).\par
По умолчанию задан МД 2053 (шестнадцатиричный код 805), канал 0, устройство 0.

Переписывается 500 зон в 300 зону системного МД.

\par\hangindent=1.5em 5) Вызвать операционную систему. Для этого необходимо выполнить последовательность действий в рамках ППО:\par
\par\hangindent=1.5em - У0\par
Для ПМ типа ДВК:

\par\hangindent=1.5em - Вызос р --- Вызов с росписью каталогов\par
Ответ системы: есть

\par\hangindent=1.5em - Ячейка 2:= 5\par
\par\hangindent=1.5em - Ячейка 3:= N пакета в шестнадцатиричном виде\par
\par\hangindent=1.5em - Ячейка а:= 0 --- Соответствует тумблерному регистру\par
\par\hangindent=1.5em - Ячейка в:= N устройства МД в шестнадцатиричном виде\par
\par\hangindent=1.5em - Ячейка с:= N канала МД в шестнадцатиричном виде\par
\par\hangindent=1.5em - У0\par
\par\hangindent=1.5em - Ар\par
\par\hangindent=1.5em - Убедиться, что РК=0 3F 07 000 (стоп 0)\par
\par\hangindent=1.5em - Директива «П»- пуск или нажатие соответствующей клавиши на ПМ;\par
Для ПМ типа ЕС-1840:

\par\hangindent=1.5em - ОС р У<N устройства с 2053> М<N машины> В<N варианта>\par
Ответ системы: начальный адрес и длина программы в словах.

По умолчанию задан МД 2053 (шестнадцатиричный код 805), канал 0, устройство 0, N МАШИНЫ=1, N ВАРИАНТА=1. Значение тумблерного регистра (ячейка «А») стандартно равно нулю (см. «ОС Диспак. Руководство системного программиста» ИЫ.00110-01 32 01).

После нормального завершения вызова ОС на операторский терминал выдаются сообщения следующего вида

<Дата> <время>

Рес: <значение обобщенного ресурса системы>

\par\hangindent=1.5em 6) После получения указанных сообщений с терминала оператора необходимо задать следующие директивы:\par
\par\hangindent=1.5em - Сме <N смены> ответ ОС: есть\par
\par\hangindent=1.5em - Дат <число> <месяц> ответ ОС: дата: <число>.<месяц>.<год>\par
\par\hangindent=1.5em - Нет арх ответ ОС: есть\par
\par\hangindent=1.5em - Вре <время> ответ ОС: дата: <число>.<месяц>.<год> время: <время> где <время> = ччммсс; ЧЧ --- Часы; ММ --- Минуты; СС --- Секунды.\par
После получения последнего сообщения можно начинать работу с ОС.

\par\hangindent=1.5em 7) В случае отсутствия любого из указанных ответов системы необходимо повторить вызов ОС с пункта 5).\par
\section*{3.~Перепись систем программирования}\addcontentsline{toc}{section}{3.~Перепись систем программирования}
\par\hangindent=1.5em 1) На работающей операционной системе директивами оператора «Есм <нбоб> <НУ>» (где <НБОБ>=2,3 --- Номера МЛ; <НУ>- логический номер магнитофона), присвоить лентам номера.\par
\par\hangindent=1.5em 2) Переписать с МЛ2 (системы программирования режима совместимости) программой «UNICOP» из файла 1 с 0 зоны в 0 зону МД 2048 2000 зон;\par
\par\hangindent=1.5em 3) Переписать с МЛ3 (система программирования) программой «UNICOP» из файла 1 с 0 зоны в 2000 зону МД 2048 2000 зон.\par
\par\hangindent=1.5em 4) После окончания переписи на операторском терминале следует набрать директивы:\par
«Дай арх»

Ответ системы: «Есть»

«Вре» ответ ОС: дата: <число>.<месяц>.<год> время: <время> где <время> = ччммсс; ЧЧ --- Часы; ММ --- Минуты; СС --- Секунды.

Убедиться в формировании задачи архива: на операторском терминале появится сообщение

Ф<NCKB> 41991981062N, где N --- Номер ЭВМ.

После формирования задачи архива система работает в стандартном режиме. На дисках 2053 и 2048 получена рабочая копия ОСПО.

\begin{center}Структура диалога\end{center}
\begin{center}\small
\begin{BVerbatim}
           -----------------
     Ответ системы         директива пользователя
     -------------         ----------------------
                               UNICOP 419955
 <<=UNICOP-NET=(ВЕРСИЯ)=>>
 РАБОТА(ЕСБ6, Б6ЕС,...):          ЕСБ6
 Том ЕС:                           2 (3 - Для МЛ3)
 Том Б6:                          2048
 Номер ФАЙЛА(0-С места):           1
 Номер зоны Б6:                    0 (2000 - Для МЛ3)
 Номер зоны ЕС:                    0
 Перемотка по окончании(да, нет):   н
 Число файлов ЕС:                  2
 Число зон:                        2000
<Для контроля повторяется набранная информация>
 О’к                           <пустая посылка>
 Терминал освободить              н
<Информация о переписанных зонах>
<Контрольная сумма переписи>
                                   КНЦ
 Конец сеанса
                                       Приложение 2.
                    Работа с документацией
\end{BVerbatim}
\end{center}

Получение документации на бумажном носителе, а также просмотр документа на терминале, осуществляется на работающей операционной системе с помощью программы «Димип». Программа обращается к области «ДОКЭБ», в которой собрана документация в виде архива программы «Димип». Структура диалога для получения и просмотра документации:

\begin{center}Структура диалога\end{center}
\begin{center}\small
\begin{BVerbatim}
         -----------------
   Ответ системы         директива пользователя
   -------------         ----------------------
                         Димип 419999
=                        КТ ДОКЭБ 0 ДОКЭБ
=                        СФ
<Список файлов>          ред <имя файла>
*                        л - Просмотр текста на терминале
*                        ф - Формирование задачи выдачи
                             Документа
*                        вый
Конец сеанса
\end{BVerbatim}
\end{center}

Ниже приведены децимальные номера документов и соответствующие им имена файлов для печати и просмотра указанных документов: децимальный номер документа файл печати файл просмотра

\begin{center}\small
\begin{BVerbatim}
---------------------------  -----------  --------------
  ИЫ.00110-01 32 01           РУСПР1           руспр
  ИЫ.00110-01 32 02           рпсовм           рпсовм
  ИЫ.00110-01 33 01           РПР1             РПР1
                              РПР21            РПР2
                              РПР3             РПР3
  ИЫ.00110-01 33 02           печасс           рпрсп
  ИЫ.00110-01 33 03           ПРЯЗА1           пряза
  ИЫ.00110-01 34 01           ро               ро
                              С                с
  ИЫ.00110-01 35 01           ОПИЯЗ1           опияз
                Лист регистрации изменений
\end{BVerbatim}
\end{center}

.-----------------------------------------------------------------.

\begin{center}\small\begin{tabularx}{\linewidth}{|>{\raggedright\arraybackslash\hsize=0.176\hsize}X|>{\raggedright\arraybackslash\hsize=0.706\hsize}X|>{\raggedright\arraybackslash\hsize=1.059\hsize}X|>{\raggedright\arraybackslash\hsize=1.059\hsize}X|>{\raggedright\arraybackslash\hsize=1.059\hsize}X|>{\raggedright\arraybackslash\hsize=1.059\hsize}X|>{\raggedright\arraybackslash\hsize=1.235\hsize}X|>{\raggedright\arraybackslash\hsize=1.588\hsize}X|>{\raggedright\arraybackslash\hsize=1.941\hsize}X|>{\raggedright\arraybackslash\hsize=1.235\hsize}X|>{\raggedright\arraybackslash\hsize=0.706\hsize}X|>{\raggedright\arraybackslash\hsize=0.176\hsize}X|}\hline
 & \small\bfseries\hspace{0pt}И з м. & \small\bfseries\hspace{0pt}Номер измененных & \small\bfseries\hspace{0pt}Лист замененных & \small\bfseries\hspace{0pt}В (СТ новых & \small\bfseries\hspace{0pt}Аниц) аннулированных & \small\bfseries\hspace{0pt}Всего листов (СТР.) В документе & \small\bfseries\hspace{0pt}Номер документа & \small\bfseries\hspace{0pt}Входящий номер сопроводи-тельного документа & \small\bfseries\hspace{0pt}Подпись & \small\bfseries\hspace{0pt}Д а т а &  \\ \hline
\end{tabularx}\end{center}

\end{document}
